\documentclass[11pt]{article}
\usepackage[margin=1.05in]{geometry}
\usepackage{amsmath,amssymb,amsthm,mathtools}
\usepackage{microtype}
\usepackage{booktabs}
\usepackage{array}
\usepackage{enumitem}
\usepackage{xcolor}
\usepackage[colorlinks=true,linkcolor=blue!50!black,citecolor=blue!50!black,urlcolor=blue!50!black]{hyperref}

\newtheorem{theorem}{Theorem}
\newtheorem{proposition}[theorem]{Proposition}
\newtheorem{lemma}[theorem]{Lemma}
\newtheorem{corollary}[theorem]{Corollary}
\newtheorem{conjecture}[theorem]{Conjecture}
\theoremstyle{definition}
\newtheorem{definition}[theorem]{Definition}
\newtheorem{problem}{Open problem}
\theoremstyle{remark}
\newtheorem{remark}[theorem]{Remark}

\newcommand{\Q}{\mathbb{Q}}
\newcommand{\B}{\mathcal{B}}
\newcommand{\rank}{\operatorname{rank}}
\newcommand{\dist}{\operatorname{dist}}
\newcommand{\Bin}{\operatorname{Bin}}
\newcommand{\CP}{\mathrm{CP}}

\title{Measurement designs and the square root\\in the subexponential \#SAT algorithm}
\author{Yamba AI Math Lab\thanks{Draft for review, prepared with substantial assistance from Claude (Anthropic) in Claude Code under the direction of Christopher Yamba. Not submitted. All claims carry the status recorded in Section~\ref{sec:status}.}}
\date{Draft of 29 October 2026, revised after a separate review}

\begin{document}
\maketitle

\begin{abstract}
A manuscript released on 26 October 2026 in the openai/math collection claims that the Exponential Time Hypothesis is false: it counts the satisfying assignments of a $k$-CNF on $n$ variables in time $2^{n/(\log n)^{1/2+o(1)}}$. The square root in the exponent enters through one step of the method, a \emph{measurement design} that separates a deformed polynomial into small independent groups. We study that step on its own. First, every design of the manuscript's own template, in which the replacement rows are built from a matrix supported on Hamming-close label pairs, needs a fraction $e^{-\Theta(R^2/Q)}$ of measurements, where $R$ is the deformation degree and $Q$ the number of interpolation nodes; the manuscript's matrix is optimal within the template up to a constant factor in $R$. Second, designs outside the template, with auxiliary matrices in the manuscript's algebra of commuting swaps and measurement matrices that may or may not depend on the node, are equivalent to a problem on the cube: attach a polynomial $p_v$ of degree $<R$ with $p_v(0)=1$ to every vertex $v$ of $\{0,1\}^Q$; the measurements needed at node $t$ are the edges $\{v,v+e_t\}$ with $p_v(x_t)+p_{v+e_t}(x_t)\ne0$. A set of edges is the bad set of a family exactly when it supports a solution of one explicit linear identity per $R$-dimensional subcube. Third, a \emph{slab} family has a fraction at most $1/R$ of bad edges in every direction for every $Q$, and needs a single measurement matrix of relative size at most $2/R$; for $Q\ge\frac98(R+2)^2$ and large $R$ it beats every template design by a factor of order $R$. A family with bad fraction $R^{-C}$, $C>3$, at $Q\approx R^{(C+1)/2}$ would improve the exponent of the method to $n(\log n)^{-(C-1)/(C+1)+o(1)}$, and such a family would have to leave almost every vertex of the cube free of bad edges. Fourth, a general rank argument requires only $4^{-R}/R$, and the natural combinatorial relaxations are exponentially weak. The core statements are verified in Lean~4 with Mathlib. All results concern the method and hold whether or not the manuscript's main theorem is correct; their novelty has not been verified.
\end{abstract}

\section{Introduction}

The Exponential Time Hypothesis (ETH) of Impagliazzo and Paturi asserts that 3-SAT on $n$ variables needs time $2^{\delta n}$ for some $\delta>0$. The openai/math manuscript \emph{The Exponential Time Hypothesis is False} (family 378, revised 25 October 2026, \cite{oai378}) claims a deterministic algorithm that counts the satisfying assignments of a $k$-CNF with $n$ variables and $m$ clauses in
\[
(n+m)^{O_k(1)}\,2^{O_k(n\log\log n/\sqrt{\log n})}
\]
time and polynomial space, and $2^{O(n/\sqrt{\log n})}$ at bounded density. An upstream Lean development proves the slightly weaker exponent $n(\log\log n)^2/\sqrt{\log n}$ for a multi-tape machine \cite{oai378lean}. An independent working group led by Impagliazzo, Paturi and Williams is examining the proof. Expositions have appeared \cite{odonnell,kintali}.

Whatever the outcome of that review, the method raises a sharp question: \emph{where does the square root come from, and can it be removed?} The algorithm deforms a polynomial whose value at $0$ is the count, recovers the deformed polynomial from Taylor data at $Q$ nodes, and evaluates each Taylor coefficient by separating the variables into $r$ groups after fixing a small number of integer \emph{measurements}. The degree $R$ of the deformation controls the cost of the Taylor data, and the number of measurements controls the cost of the separation. In the manuscript's construction the measurement fraction decays like $e^{-R^2/(8Q)}$, so $R$ must be of order $\sqrt{Q\log(1/\varepsilon)}$, and the square root follows.

\paragraph{Results.} We isolate the measurement step and prove:
\begin{enumerate}[leftmargin=2em]
\item \textbf{A barrier for the manuscript's template} (Theorem~\ref{thm:barrier}, Corollary~\ref{cor:degree}). Every pair $(M_0,C)$ satisfying the manuscript's two conditions (2.14) needs $c/D\ge\Pr[X\ge k+R]\Pr[X\le k]$ with $X$ binomial, hence $R\gtrsim\sqrt{Q\log(1/\varepsilon)}$. The manuscript's matrix $V_0f(W_0)$ is optimal within the template up to a constant factor in $R$.
\item \textbf{A reduction for designs beyond the template} (Proposition~\ref{prop:reduction}). Inside the commuting swap algebra, designs are exactly families of polynomials on the cube, and the measurements at node $t$ are the weighted bad edges in direction $t$. We call the resulting problem the \emph{cube polynomial problem} $\CP(Q,R)$.
\item \textbf{A local linear form} (Theorem~\ref{thm:local}). An edge function is the defect function of a family if and only if it satisfies one explicit identity per $R$-dimensional subcube; so $\CP(Q,R)$ is a sparse-solution problem for an explicit linear system, and every obstruction is a combination of subcube obstructions.
\item \textbf{Global formulas and the slab family} (Lemma~\ref{lem:global}, Theorem~\ref{thm:slab}, Corollary~\ref{cor:helps}). There are explicit families with bad fraction at most $1/R$ in every direction for every $Q$, which need a single measurement matrix of relative size at most $2/R$; for $Q\ge\frac98(R+2)^2$ and large $R$ they beat every template design by a factor of order $R$.
\item \textbf{The target, and what it forces} (Propositions~\ref{prop:criterion} and~\ref{prop:forces}). Bad fraction $R^{-C}$ with $C>3$ at $Q\approx R^{(C+1)/2}$ would improve the exponent of the method ($C>4$ under the manuscript's own row count), and any design with exponent below $n/\sqrt{\log n}$ must leave almost every vertex of the cube free of bad edges. A lower bound $c/R$ on the measure of the touched vertices, which the slab attains, would therefore prove a square-root barrier for all swap designs (Corollary~\ref{cor:criterion}).
\item \textbf{Barriers and duality} (Theorems~\ref{thm:general}, \ref{thm:stress}, \ref{thm:product}). Every design, in or out of the swap algebra, needs only $c_t/D\ge4^{-R}/R$ from a forced-identity argument; a set of good edges is feasible if and only if every stress on it has total weight zero; obstructions have at least $2^R$ vertices, so fractional and subcube-hitting relaxations give only exponentially small bounds; products of families preserve the bad fraction.
\end{enumerate}
Small optima, computed by integer programming, re-verified in rational arithmetic and in five cases certified exactly, beat the slab in every case searched with $3\le R\le Q-1$ (Section~\ref{sec:data}). The decisive asymptotic question, whether the bad fraction can be made polynomially small in $R$ at $Q\gg R^2$, is open (Section~\ref{sec:open}).

\paragraph{Status.} The statements concern the method, not the truth of the main theorem: they hold whether or not the manuscript's algorithm is correct. The core lemmas are machine-checked (Section~\ref{sec:lean}). A first draft attributed the advantage of the slab family to node-dependent measurements; a separate review showed that a single measurement matrix already achieves it, and that the dividing line is the manuscript's template. Our literature check (Section~\ref{sec:status}) found no prior statement of these results, but it was not a systematic review, so their novelty is \textbf{not verified}.

\section{The measurement step of the algorithm}\label{sec:setting}

We follow the bounded-occurrence overview of \cite[Section~2]{oai378}, with its numbering; all costs ignore polynomial factors. A formula with $n$ variables and $M\le3n$ literal occurrences gives $\#F=\sum_{x\in\{0,1\}^n}\prod_e\Phi_e(x)$. Each occurrence of $X_v$ in clause $e$ is replaced by a linear form $\langle\vec P_{e,v}(T),\vec X\rangle$ whose coefficients are polynomials of degree at most $R$ in a deformation variable $T$, with $\vec P_{e,v}(0)=\delta_v$. The resulting $H(T)$ satisfies $H(0)=\#F$ and is recovered by Hermite interpolation from its Taylor coefficients at nodes $t=1,\dots,Q$; each Taylor coefficient is assembled from Boolean substitutions by inclusion and exclusion (Lemma~3.3). At each node the deformed clauses are separated: an $(r,\tau,D)$-design at $t$ (Definition~2.1; in general form Definition~3.1, with groups and matrices $C_{t,g}$ for every node $t$ and group $g$) consists of $r$ groups of at most $2n/r$ variables and measurement matrices with at most $2\tau n$ rows in total, each of $\ell_1$ norm at most $D$, such that the dependence of each replacement form on any group other than its clause's own is a combination of that group's measurements. The three costs multiply to
\begin{equation}\label{eq:cost}
\underbrace{\sum_{j\le q}2^j\binom nj}_{\text{Taylor data},\ q=\lfloor3nR/Q\rfloor}\cdot\underbrace{(2D+1)^{2\tau n}}_{\text{measurements}}\cdot\underbrace{2^{2n/r}}_{\text{one group}}
\end{equation}
\cite[display (2.11)]{oai378}, which is $2^{O(n[(R/Q)\log(2Q/R)+\tau\log(2D+1)+1/r])}$.

\paragraph{The manuscript's template.} The designs are built on auxiliary labels $u=(a_u,d_u)\in\B=[r]^Q\times[r]^Q$, so $D=|\B|=r^{2Q}$; for labels $u,w$ put $E_{u,w}=\{t:a_w(t)\ne d_u(t)\}$. The construction takes a matrix $M_0\in\Q^{\B\times\B}$ and a measurement matrix $C\in\Q^{c\times\B}$ with
\begin{enumerate}[label=(\roman*),leftmargin=2em]
\item (kernel) $M_0\xi=\xi$ for every $\xi\in\ker C$, and
\item (support) $M_0[u,w]\ne0\Rightarrow|E_{u,w}|\le R-1$,
\end{enumerate}
which are the conditions (2.14), builds the replacement rows (2.16) and (4.10) from them, and requires $c\le D\tau/(3r)$. We call designs of this form the \emph{template}. The manuscript uses the commuting swaps $\Pi_s$ that exchange the two coordinates at position $s$, the complete swap $V_0=\prod_s\Pi_s$, the operator $W_0=\frac12\sum_s(I-\Pi_s)$, and $M_0=V_0f(W_0)$, where $f$ has degree $R-1$ and interpolates $(-1)^j$ on a window $\mathcal I$ of $R$ consecutive integers around $Qp$, $p=\frac{r-1}{2r}$. A $\{0,\pm1\}$ eigenbasis gives a matrix $C$ with $c\le D\Pr[\Bin(Q,p)\notin\mathcal I]\le2De^{-R^2/(8Q)}$ rows (Lemma~4.1 and Step~5). The constraints (2.17) include $12D^2\le M$ and $2Qr^2\le n$.

\paragraph{Counting rows.} Definition~3.1 counts measurement rows after the split into groups, $\sum_gc_{t,g}\le2\tau n$. The manuscript's own construction uses $r$ rows per measurement, zero rows included; a design may drop zero rows. Both counts appear below; we call them \emph{the manuscript's count} and \emph{the sharper count}.

\section{Template designs: an isoperimetric barrier}\label{sec:single}

\begin{theorem}[rank barrier]\label{thm:barrier}
Let $(M_0,C)$ satisfy (i) and (ii), and let $A,B\subseteq[r]^Q$ with $\dist_H(a,b)\ge R$ for all $a\in A$, $b\in B$. Then $c\ge\rank(I-M_0)\ge|A|\,|B|$.
\end{theorem}

\begin{proof}
By (i), $\ker C\subseteq\ker(I-M_0)$, so $\rank(I-M_0)\le\rank C\le c$. Let $U=\{u:a_u\in A,\ d_u\in B\}$. For $u,w\in U$, $|E_{u,w}|=\dist_H(d_u,a_w)\ge R$, so $M_0[u,w]=0$ by (ii). The principal submatrix of $I-M_0$ on $U$ is the identity, so $\rank(I-M_0)\ge|U|=|A||B|$.
\end{proof}

\begin{corollary}[binomial form]\label{cor:binomial}
Let $h=\lfloor r/2\rfloor$ and $X\sim\Bin(Q,h/r)$. For every integer $k$, $c/D\ge\Pr[X\ge k+R]\cdot\Pr[X\le k]$.
\end{corollary}

\begin{proof}
Fix $H\subseteq[r]$ with $|H|=h$ and let $N_H(a)$ count the positions with $a(t)\in H$. Then $\dist_H(a,b)\ge N_H(a)-N_H(b)$. Apply Theorem~\ref{thm:barrier} to $A=\{N_H\ge k+R\}$ and $B=\{N_H\le k\}$ and divide by $D=r^{2Q}$.
\end{proof}

\begin{corollary}[necessary degree]\label{cor:degree}
If $c\le\varepsilon D$ then $(R+1)^2\ge Q\,p_h(1-p_h)\ln\frac1{2(Q+1)\varepsilon}$ with $p_h=\lfloor r/2\rfloor/r\in[\frac13,\frac12]$, provided $\lceil Qp_h\rceil+R\le Q$ (otherwise $R>Q/2-1$).
\end{corollary}

\begin{proof}
Take $k=\lceil Qp_h\rceil$, so $\Pr[X\le k]\ge\frac12$ since a binomial median lies in $\{\lfloor Qp_h\rfloor,\lceil Qp_h\rceil\}$. For $m=k+R$, $\Pr[X\ge m]\ge\binom Qm p_h^m(1-p_h)^{Q-m}\ge\frac1{Q+1}e^{-Q\,\mathrm{KL}(m/Q\,\|\,p_h)}$, and $\mathrm{KL}(q\|p)\le(q-p)^2/(p(1-p))$ with $m-Qp_h\le R+1$.
\end{proof}

The manuscript achieves $c\le D\tau/(3r)$ with $R\approx\sqrt{8Q\ln(6r/\tau)}$, while Corollary~\ref{cor:degree} forces $R\gtrsim\sqrt{\frac29Q\ln\frac{3r}{2(Q+1)\tau}}$. So \textbf{no pair satisfying (i) and (ii), in or out of the swap algebra, improves on $V_0f(W_0)$ by more than a constant factor in $R$}: about $6$ when $\ln(r/\tau)\gg\ln Q$, and about $\sqrt{72}\approx8.5$ in the manuscript's own regime, where $r/\tau=r^2\log_2n$ and $\ln(r/\tau)\approx2\ln Q$. Within the template the square root is forced. For the swap algebra and the uniform measure, the rank identity of O'Donnell \cite[Lemma~4.1]{odonnell} combined with the sign-error bound of Aspnes, Beigel, Furst and Rudich \cite{abfr} gives a similar conclusion; Theorem~\ref{thm:barrier} needs neither.

\section{Designs beyond the template and the cube polynomial problem}\label{sec:cp}

The design definitions ask only that, at each node, the nonlocal part of every replacement row be determined by the measurements. Theorem~\ref{thm:barrier} uses the special form of the template rows and does not apply to other designs. We now describe all designs whose auxiliary matrix lies in the swap algebra.

Each position $s$ splits $\Q^{[r]\times[r]}$ into the swap-symmetric part (dimension $r(r+1)/2$) and the antisymmetric part (dimension $r(r-1)/2$). A product eigenvector has a type $\varepsilon\in\{0,1\}^Q$ ($\varepsilon_s=1$ when antisymmetric at $s$), and $\Pi_S=\prod_{s\in S}\Pi_s$ acts on it by $\chi_S(\varepsilon)=(-1)^{|S\cap\varepsilon|}$. Under the uniform measure on labels the type has the product law $\mu_p$, $p=\frac{r-1}{2r}$.

A \emph{swap design of degree $<R$} is $P(T)=\sum_{S\subseteq[Q]}c_S(T)\Pi_S$ with $P(0)=I$ and $\deg c_S<R$. At node $t$ the row of label $u$ may read freely the labels $w$ with $a_w(t)=d_u(t)$; the rest of $P(x_t)$ is the nonlocal part $N_t$. Since $\rank N_t\le\sum_g\rank N_t^{(g)}\le\sum_gc_{t,g}$, at least $\rank N_t$ rows are needed at node $t$ after the split; conversely the $\{0,\pm1\}$ eigenvectors spanning the row space of $N_t$ suffice, each meeting at most two groups at the node. A single measurement matrix serving every node needs $\dim\sum_t\operatorname{rowspace}N_t$ rows before the split.

\begin{proposition}[reduction]\label{prop:reduction}
Put $p_\varepsilon(T)=\sum_Sc_S(T)\chi_S(\varepsilon)$ for $\varepsilon\in\{0,1\}^Q$. Then $P(0)=I$ if and only if $p_\varepsilon(0)=1$ for all $\varepsilon$; $\deg P<R$ if and only if $\deg p_\varepsilon<R$ for all $\varepsilon$; and for every node $t$
\[
\frac{\rank N_t}{D}=\frac{r-1}{r}\,\mu_p\big(\{\varepsilon':\ p_{\varepsilon'}(x_t)+p_{\varepsilon'+e_t}(x_t)\ne0\}\big),
\]
where $\varepsilon'$ ranges over types with $\varepsilon'_t=0$ and $\mu_p$ is the product measure on the other $Q-1$ coordinates.
\end{proposition}

\begin{proof}
The first two statements are Fourier inversion on $\{0,1\}^Q$. The entry $(u,w)$ of $\Pi_S$ is nonzero only for $w=\mathrm{swap}_S(u)$. If $t\in S$ then $a_w(t)=d_u(t)$ and the entry is local. If $t\notin S$ then $a_w(t)=a_u(t)$, which equals $d_u(t)$ exactly when $u$ is diagonal at $t$. Hence $N_t=D_t\sum_{S\not\ni t}c_S(x_t)\Pi_S$ with $D_t$ the projector onto labels off-diagonal at $t$. The projector commutes with every swap, has rank $r(r-1)$ at position $t$, and the remaining factor is diagonal in the product eigenbasis with eigenvalue $\frac12(p_{\varepsilon'}(x_t)+p_{\varepsilon'+e_t}(x_t))$ on the eigenvectors of type $\varepsilon'$ at the other positions.
\end{proof}

We checked the formula exactly on the full matrices, masking their nonlocal entries and computing ranks modulo two 31-bit primes, for several families in each of the six cases $(r,Q,R)=(2,3,2)$, $(2,3,3)$, $(3,2,2)$, $(2,4,2)$, $(2,4,3)$ and $(3,3,2)$.

\begin{definition}[cube polynomial problem]
Fix distinct nonzero nodes $x_1,\dots,x_Q$. A \emph{family} assigns to every vertex $v\in\{0,1\}^Q$ a polynomial $p_v$ of degree $<R$ with $p_v(0)=1$. The \emph{defect} of the edge $e=\{v,v+e_t\}$ is $d_e=p_v(x_t)+p_{v+e_t}(x_t)$, and the edge is \emph{bad} if $d_e\ne0$. The value of a family is $\beta=\max_t\beta_t$, where $\beta_t$ is the measure of the bad edges in direction $t$ (the measure of their endpoints with $v_t=0$, on the other coordinates). $\CP(Q,R)$ asks for the least value $\beta^*(Q,R)$.
\end{definition}

By Proposition~\ref{prop:reduction}, a family of value $\beta$ for $\mu_p$ is a swap design with $\rank N_t/D\le\frac{r-1}r\beta$ at every node. The manuscript's template rows are not in the swap algebra, but at every node their nonlocal parts coincide entrywise with those of the \emph{window family} $p_\varepsilon=1-\sigma_\varepsilon f(|\varepsilon|)+\sigma_\varepsilon G(T,\varepsilon)$, $\sigma_\varepsilon=(-1)^{|\varepsilon|}$, with $G$ the global formula of Lemma~\ref{lem:global} below (checked in four small cases by the separate review). Edges with both weights in the window are good, so its value is at most the binomial tail $\le2e^{-R^2/(8Q)}$.

\begin{proposition}[subcube obstruction]\label{prop:subcube}
For any $R$ distinct directions, every $R$-dimensional subcube in those directions contains a bad edge. More generally, in every $d$-dimensional subcube at most $R-1$ directions are free of bad edges. Consequently $\sum_{t\in T}\beta_t\ge\min(p,1-p)^{R-1}$ for every set $T$ of $R$ directions, and $\beta^*(Q,R)\ge2^{1-R}/R$ for the uniform measure.
\end{proposition}

\begin{proof}
On a $d$-subcube put $H=\sum_vp_v$, of degree $<R$ with $H(0)=2^d\ne0$. For a direction $t$ of the subcube, $H(x_t)$ is the sum of the defects of the $t$-edges, so it vanishes if they are all good; a nonzero polynomial of degree $<R$ has fewer than $R$ roots. For $d=R$ this is the first statement; the bound follows by summing over the bases of the subcubes.
\end{proof}

\section{A local linear form}\label{sec:local}

For a set $T$ of $R$ directions let $\lambda^T_t=\prod_{s\in T\setminus t}x_s/(x_s-x_t)$ $(t\in T)$, the Lagrange weights at $0$, so that $f(0)=\sum_{t\in T}\lambda^T_tf(x_t)$ for every $f$ of degree $<R$.

\begin{theorem}[subcube identities characterize defects]\label{thm:local}
Let $1\le R\le Q$. An edge function $d$ is the defect function of a family if and only if, for every $R$-dimensional subcube $\Gamma$ with direction set $T$,
\[
\sum_{e\in E(\Gamma)}\lambda^T_{t(e)}\,d_e\ =\ 2^R .
\]
\end{theorem}

\begin{proof}
Necessity: sum $p_v(0)=\sum_{t\in T}\lambda^T_tp_v(x_t)$ over the $2^R$ vertices of $\Gamma$ and split each inner sum into the $t$-edges. Sufficiency: let $d_t(v)=d_{\{v,v+e_t\}}$, a function on the cube that is even in coordinate $t$, and let $A_T$ average over the $T$-subcube through a vertex. The identity for the subcube through $b$ reads $\sum_{t\in T}\lambda^T_t(A_Td_t)(b)=2$. Since $A_T\chi_S=\chi_S$ if $S\cap T=\varnothing$ and $0$ otherwise, the identities say that $\sum_{t\in T}\lambda^T_t\widehat{d_t}(S)=2[S=\varnothing]$ for every $S$ and every $R$-set $T\subseteq[Q]\setminus S$. If $|[Q]\setminus S|\ge R$, values whose extrapolation to $0$ from every $R$-subset gives the same number lie on one polynomial of degree $<R$; so there are $h_S$ of degree $<R$ with $h_S(0)=2[S=\varnothing]$ and $\widehat{d_t}(S)=h_S(x_t)$ for $t\notin S$. If $|[Q]\setminus S|<R$, choose $h_S$ with $h_S(0)=0$ and these values. Then $p_v=\frac12\sum_Sh_S\chi_S(v)$ has degree $<R$, $p_v(0)=1$, and defects $\sum_{S\not\ni t}\widehat{d_t}(S)\chi_S(v)=d_t(v)$.
\end{proof}

So $\beta^*(Q,R)$ is the least worst-direction support of a solution of the subcube system $Md=2^R\mathbf1$, with one row per $R$-subcube. By the Fredholm alternative, a set $G$ of edges can be made simultaneously good if and only if no combination of subcube stresses (Theorem~\ref{thm:stress}) with nonzero total vanishes on every edge outside $G$. We checked the theorem against the polynomial formulation on 557 edge sets in exact arithmetic, including all optimal and slab sets of Table~\ref{tab:data} and one-edge modifications of them.

\section{Global formulas and the slab family}\label{sec:slab}

\begin{lemma}[global formulas]\label{lem:global}
For coefficients $\gamma_S$ $(S\subseteq[Q])$ and distinct nonzero nodes put
\[
G(T,\varepsilon)=\sum_S\gamma_S\,\chi_S(\varepsilon)\prod_{s\in S}\Big(1-\frac T{x_s}\Big).
\]
Then $G(x_t,\varepsilon)=G(x_t,\varepsilon+e_t)$ for every vertex and direction, $\deg_TG\le\max\{|S|:\gamma_S\ne0\}$, and $G(0,\varepsilon)=\sum_S\gamma_S\chi_S(\varepsilon)$.
\end{lemma}

\begin{proof}
Terms with $t\in S$ vanish at $T=x_t$, and terms with $t\notin S$ satisfy $\chi_S(\varepsilon+e_t)=\chi_S(\varepsilon)$.
\end{proof}

So $p_\varepsilon=\sigma_\varepsilon G(\cdot,\varepsilon)$ makes every edge good, and the only defect is the normalization $p_\varepsilon(0)=\sigma_\varepsilon g(\varepsilon)$ with $g=\sum_S\gamma_S\chi_S$ of Fourier degree $<R$. Since parity has Fourier degree $Q$, no all-good family exists for $Q\ge R$.

\begin{lemma}[symmetric functions have low degree]\label{lem:symmetric}
Over a field of characteristic $0$, if $\deg f\le d$ then $\varepsilon\mapsto f(|\varepsilon|)$ is a combination of the characters $\chi_S$ with $|S|\le d$.
\end{lemma}

\begin{proof}
$|\varepsilon|=\sum_s\frac{1-\chi_{\{s\}}(\varepsilon)}2$, and multiplication by $\chi_{\{s\}}$ maps $\chi_S$ to $\chi_{S\triangle\{s\}}$, raising the degree by at most one.
\end{proof}

\begin{theorem}[slab family]\label{thm:slab}
Fix $R\ge1$ and an offset $s_0$. Cut the weights $0,\dots,Q$ into slabs $\{w:\lfloor(w+s_0)/R\rfloor=k\}$ of at most $R$ consecutive values. On slab $k$ let $f_k$ interpolate $(-1)^w$ (degree $<R$), let $\gamma^{(k)}$ be the Fourier coefficients of $f_k(|\cdot|)$ given by Lemma~\ref{lem:symmetric}, and put $p_\varepsilon=\sigma_\varepsilon G_{\gamma^{(k)}}(\cdot,\varepsilon)$ for $|\varepsilon|$ in slab $k$. Then $\deg p_\varepsilon<R$, $p_\varepsilon(0)=1$, every edge inside a slab is good, and for some $s_0<R$, for every direction $t$ and every $Q\ge1$, at most $2^{Q-1}/R$ edges in direction $t$ are bad. Averaging over $s_0$ gives $\mu_p(\text{bad in direction }t)\le1/R$ for every product measure with identical coordinate laws.
\end{theorem}

\begin{proof}
Degree and normalization follow from Lemmas~\ref{lem:global} and~\ref{lem:symmetric}. An edge in direction $t$ joining weights $w$ and $w+1$ is good when both lie in one slab. Each $w$ is the top of a slab for exactly one offset modulo $R$, so summing over the $R$ offsets counts every edge at most once; the count does not depend on $t$.
\end{proof}

\begin{corollary}[leaving the template helps]\label{cor:helps}
The slab design needs $\rank N_t/D\le\frac{r-1}{rR}$ at every node, and, for a suitable offset, a single measurement matrix with $c/D\le2/R$, for every $Q$. If $Q\ge\frac98(R+2)^2$ and $Q$ exceeds an absolute constant, every template design has $c/D\ge0.01$.
\end{corollary}

\begin{proof}
The first bound is Proposition~\ref{prop:reduction}. The bad edges join weights $m$ and $m+1$ with $m+1+s_0\equiv0\pmod R$, so the nonlocal rows lie in the span of the eigenvectors whose type has weight $\equiv-s_0-1$ or $-s_0\pmod R$; a single matrix spanning these serves every node, and averaging over $s_0$ gives $2/R$. For the template, take $k=\lceil Qp_h\rceil$ in Corollary~\ref{cor:binomial}, so $\Pr[X\le k]\ge\frac12$; with $\sigma=\sqrt{Qp_h(1-p_h)}\ge\sqrt{2Q/9}$ the hypothesis gives $k+R\le Qp_h+2\sigma$, and the Berry and Esseen theorem gives $\Pr[X\ge Qp_h+2\sigma]\ge0.02$ for large $Q$.
\end{proof}

So in that regime leaving the template wins by a factor of order $R$, and node dependence matters only by a constant factor for the slab.

\begin{proposition}[improvement criterion]\label{prop:criterion}
Assume the manuscript's cost analysis \eqref{eq:cost}, its constraints (2.17) and the requirement $c_t/D\le\tau/(3r)$ at every node, applied with swap designs. Suppose that for some $C$ there are families with $\beta\le R^{-C}$ (for $\mu_p$) at $Q\approx R^\alpha$ for all large $R$. With zero rows dropped, the exponent of the method becomes $n(\log n)^{-\min(\alpha-1,\,C-\alpha)/\alpha+o(1)}$, best at $\alpha=(C+1)/2$, which beats $n/(\log n)^{1/2}$ exactly when $C>3$. Under the manuscript's count it is $n(\log n)^{-\min(\alpha-1,\,(C-\alpha)/2)/\alpha+o(1)}$, best at $\alpha=(C+2)/3$, which beats the square root exactly when $C>4$.
\end{proposition}

\begin{proof}[Derivation]
Write $\lambda=\log_2n$, $Q\approx\lambda/(4\log_2r)$ and $R=\lambda^b$. The Taylor term is about $\lambda^{b-1}$ up to logarithms. With zero rows dropped, $\tau\ge6\beta$ suffices, the measurement term is about $6\lambda\beta=6\lambda^{1-bC}$, independent of $r$, and $r$ can be taken polynomially large, so the exponent is $\lambda^{-\min(1-b,\,bC-1)}$. Under the manuscript's count, $\tau\ge3(r-1)\beta$, the measurement term is about $3r\lambda\beta$, and balancing it with the group term $2/r$ at $r=\lambda^a$ gives $a\le\min(1-b,(bC-1)/2)$. Put $b=1/\alpha$.
\end{proof}

The slab has $C=1$; the window family is excellent at fixed ratio $Q/R$ but needs $R\gtrsim\sqrt{Q\log Q}$. If $\beta$ could be made $2^{-\Theta(R)}$, the same computation with $R=\Theta(\log\log n)$ would give $n(\log\log n)^{O(1)}/\log n$.

\section{What an improvement forces}\label{sec:forces}

For a family, the \emph{touched set} $W$ is the set of vertices incident to a bad edge. Every bad edge has both endpoints in $W$, and the endpoints of the bad edges in direction $t$ have $\mu_p$-measure exactly $\beta_t$, so $\mu_p(W)\le\sum_t\beta_t\le Q\beta$.

\begin{proposition}[what exponent $e$ needs]\label{prop:forces}
In the setting of Proposition~\ref{prop:criterion}, a design with exponent $e$ per variable has $R\le(e+1/n)Q/3$, $Q<\lambda/4$, and $\mu_p(W)\le Q\beta\le e/24$ with zero rows dropped ($\le e^2/12$ under the manuscript's count).
\end{proposition}

\begin{proof}
The Taylor term is at least $2^q$ with $q=\lfloor3nR/Q\rfloor$, so $e\ge3R/Q-1/n$. The constraint $12D^2\le M\le3n$ gives $4Q\log_2r\le\lambda-2$. The measurement term is at least $4\tau Q\log_2r$, which is at least $24Q\beta$ with zero rows dropped; under the manuscript's count it is at least $12(r-1)\beta Q\ge12Q\beta/e$, since the group term gives $r-1\ge r/2\ge1/e$.
\end{proof}

So an exponent $\lambda^{-\alpha}$ with $\alpha>\frac12$ needs $R<\sqrt Q$ up to constants and leaves almost every vertex untouched. The slab touches about $2/R$ of the vertices; the window family touches few, but needs $R\gtrsim\sqrt{Q\log Q}$.

\begin{corollary}[a barrier criterion]\label{cor:criterion}
Suppose that every family with $R\le c_0\sqrt Q$ has $\mu_p(W)\ge c/R$ for all $p\in[\frac14,\frac12]$. Then every swap design has exponent $e\ge\min(6c_0,\sqrt{288c})\,\lambda^{-1/2}$ with zero rows dropped. Under the manuscript's count the weaker hypothesis $\mu_p(W)\ge c/R^2$ already gives $e\ge\min(6c_0,(1728c)^{1/4})\,\lambda^{-1/2}$.
\end{corollary}

\begin{proof}
Let $\kappa$ be the stated minimum and suppose $e<\kappa\lambda^{-1/2}$. Since $\lambda\ge4Q$, $e<\kappa/(2\sqrt Q)$ and $R\le eQ/3<\kappa\sqrt Q/6\le c_0\sqrt Q$, so the hypothesis applies. With zero rows dropped, $c/R\le e/24$, so $eR\ge24c$, while $eR\le e^2Q/3<\kappa^2/12\le24c$. Under the manuscript's count, $c/R^2\le e^2/12$, so $e^2R^2\ge12c$, while $e^2R^2<\kappa^4/144\le12c$.
\end{proof}

Equivalently, up to a factor $Q$, $\CP(Q,R)$ asks for the sparsest vertex set $U$ whose deletion leaves an induced subgraph on which one family is entirely good: the touched set of a family is such a set, and a family on the rest extended by $p_u=1$ has $\beta\le\mu_p(U)/\min(p,1-p)$. Every such $U$ meets every $R$-subcube, but this combinatorial condition allows density $(Q+1)2^{-R/2+o(R)}$ (Section~\ref{sec:barriers}).

\section{Barriers and duality}\label{sec:barriers}

\begin{theorem}[general rank barrier]\label{thm:general}
Let $P(T)$ be any $\B\times\B$ matrix of polynomials with $\deg P<R$ and $P(0)=I$, not necessarily in the swap algebra. Let $x_1,\dots,x_R$ be distinct nodes, node $i$ using the groups of position $\pi(i)$, with $\pi$ injective, and split $P(x_i)=L_i+N_i$ where $L_i[u,w]\ne0$ implies $a_w(\pi(i))=d_u(\pi(i))$. For product sets $A=\prod_tS^A_t$ and $B=\prod_tS^B_t$ with $S^A_{\pi(i)}\cap S^B_{\pi(i)}=\varnothing$ for all $i$,
\[
|A|\,|B|\le\sum_{i=1}^R\rank N_i,\qquad\text{hence}\qquad\max_i\frac{\rank N_i}{D}\ge\frac1R\Big(\frac{h(r-h)}{r^2}\Big)^R ,
\]
which is $4^{-R}/R$ for even $r$ and $h=r/2$.
\end{theorem}

\begin{proof}
Lagrange extrapolation gives $I=P(0)=\sum_i\lambda_iP(x_i)$, hence $\sum_i\lambda_iN_i=I-\sum_i\lambda_iL_i$. On $U=\{u:a_u\in A,\ d_u\in B\}$ every $L_i$ vanishes, so the principal submatrix of $\sum_i\lambda_iN_i$ on $U$ is the identity and $|U|\le\sum_i\rank N_i$.
\end{proof}

For this forced-identity argument product sets are optimal, so it cannot give more than $4^{-R}/R$, and plugging that into Proposition~\ref{prop:criterion} shows it cannot force more than $R=\Theta(\log\log n)$: the square root is a property of the template, not of the measurement framework.

\begin{theorem}[stress duality]\label{thm:stress}
Put $\xi_t=(x_t,x_t^2,\dots,x_t^{R-1})$. For a set $G$ of edges, a \emph{stress} is a weighting $w:G\to\Q$ with $\sum_{t:\,e(v,t)\in G}w_{e(v,t)}\xi_t=0$ at every vertex $v$. There is a family in which every edge of $G$ is good if and only if every stress on $G$ has total weight $\sum_ew_e=0$. The $R$-subcube in directions $T$ carries the stress $w_e=\lambda^T_{t(e)}$, of total $2^{R-1}$.
\end{theorem}

\begin{proof}
If a family exists, vanishing moments give $\sum_tw_{e(v,t)}p_v(x_t)=(\sum_tw_{e(v,t)})p_v(0)$ at each vertex; summing over vertices, the right side is $2\sum_ew_e$ and the left side is $\sum_ew_ed_e=0$. Conversely, by linear duality the system in the coefficients of the $p_v$ is unsolvable exactly when some combination of its equations has zero left side and nonzero right side, and the multipliers of such a combination form a stress with nonzero total.
\end{proof}

\begin{corollary}[relaxations are weak]\label{cor:weak}
Every stress with nonzero total is supported on at least $2^R$ vertices. Consequently the fractional relaxation of $\CP(Q,R)$ over all obstructions has value at most $2^{1-R}/R$, and the integral subcube-hitting relaxation has value at most $2^{-R/2+o(R)}$ by the constructions of Alon \cite{alon} and of Ellis, Ivan and Leader \cite{eil}. A polynomial lower bound $\beta^*\ge R^{-O(1)}$ must exploit integrality for the algebraic obstructions, which by Theorem~\ref{thm:local} are combinations of subcube stresses.
\end{corollary}

\begin{proof}
Any $R-1$ of the vectors $\xi_t$ are independent (Vandermonde), so every vertex touched by a nonzero stress carries at least $R$ weighted edges, and a subgraph of the cube of minimum degree at least $R$ has at least $2^R$ vertices. The fractional bound follows by giving weight $1/(R2^{R-1})$ to every edge, and the hitting bound by taking all edges at a vertex set of density $2^{-R/2+o(R)}$ that meets every $R$-subcube.
\end{proof}

\begin{theorem}[products]\label{thm:product}
If $(p_v)$ is a family on $Q_1$ directions of degree $<R_1$ and $(s_u)$ a family on $Q_2$ other directions of degree $<R_2$, with all nodes distinct, then $p_{(v,u)}=p_vs_u$ is a family on $Q_1+Q_2$ directions of degree $<R_1+R_2-1$ whose bad edges are contained in those of its factors. Hence $\beta^*(Q_1+Q_2,R_1+R_2-1)\le\max(\beta^*(Q_1,R_1),\beta^*(Q_2,R_2))$ for the node blocks used.
\end{theorem}

\begin{proof}
If $p(y)+p'(y)=0$ then $(ps)(y)+(p's)(y)=0$.
\end{proof}

Products preserve the ratio $Q/R$ and the bad fraction, so they cannot reach the regime $Q\gg R^2$ that Proposition~\ref{prop:criterion} needs; they do turn small optima into infinite families. Scaling all nodes by a common factor is a symmetry, so the optimum at $(4,3)$ (nodes $t/4$) gives $\beta\le1/8$ at $(Q,R)=(4m,2m+1)$ for every $m$, using the disjoint blocks $\{5^jt/4\}$.

\section{Small cases}\label{sec:data}

We searched $\CP(Q,R)$ for the nodes $x_t=t/Q$ and the uniform measure with a mixed-integer program (HiGHS through SciPy, 240 seconds per instance, polynomial coefficients boxed by $|a|\le200$), re-checking every reported good-edge set in exact rational arithmetic, so every entry of Table~\ref{tab:data} is an upper bound on the optimum. Solver optimality and the bracketed solver bounds hold only within the coefficient box; entries marked $\dagger$ are certified optimal with no box by exhaustive exact search over hitting sets. The window column is the actual window family at its best position.

\begin{table}[h]
\centering\small
\begin{tabular}{rrllll}
\toprule
$Q$ & $R$ & subcube hitting & $\CP$ & slab & window family\\
 & & total / max & total / max & total / max & total / max\\
\midrule
3 & 2 & 3 / 1 & 4 / 2 $\dagger$ & 6 / 2 & 6 / 2\\
4 & 2 & 8 / 2 & 12 / 4 & 16 / 4 & 16 / 4\\
4 & 3 & 3 / 1 & 4 / 1 $\dagger$ & 8 / 2 & 8 / 2\\
5 & 2 & 24 / 5 & $\le32\ [25]$ / $\le7\ [5]$ & 40 / 8 & 40 / 8\\
5 & 3 & 8 / 2 & $\le14\ [8]$ / $\le4\ [2]$ & 25 / 5 & 30 / 6\\
5 & 4 & 3 / 1 & 4 / 1 & 10 / 2 & 10 / 2\\
6 & 3 & $\le22$ / 4 & $\le60\ [9]$ / $\le12\ [2]$ & 60 / 10 & 72 / 12\\
6 & 4 & 8 / 2 & $\le18\ [4]$ / $\le5\ [1]$ & 36 / 6 & 42 / 7\\
6 & 5 & 3 / 1 & 4 / 1 & 12 / 2 & 12 / 2\\
\bottomrule
\end{tabular}
\caption{Bad edges in small cases of $\CP(Q,R)$ (nodes $t/Q$, uniform measure). Also certified: $(2,2)$, $(3,3)$ and $(4,4)$ have a single bad edge.}\label{tab:data}
\end{table}

Observations, all empirical: realizability is a real constraint (certified at $(3,2)$ and $(4,3)$, where the optimum exceeds the hitting value); for $R=Q-1$ the families found have four bad edges in total; and the slab is beaten in every case searched with $3\le R\le Q-1$, and at $(5,2)$, where a family with at most $7$ bad edges in every direction exists. The families found are not unions of global formulas on level sets. Iteratively reweighted $\ell_1$ minimization on the subcube system of Theorem~\ref{thm:local} finds sparse but direction-unbalanced solutions and does not beat the slab's worst direction for $Q\ge6$, so convex relaxations are not tight here.

\section{Open problems}\label{sec:open}

\begin{problem}[the decisive question]
Determine the growth of $\beta^*(Q,R)$ for $Q\ge R^{2+\delta}$. A family with $\beta\le R^{-C}$, $C>3$, at $Q\approx R^{(C+1)/2}$ improves the exponent of the method (Proposition~\ref{prop:criterion}); a proof that $\beta^*(Q,R)\ge R^{-3-o(1)}$ there, or the touched-set bound of Corollary~\ref{cor:criterion}, shows that no swap design beats the square root.
\end{problem}

\begin{problem}[the infinite-ratio limit]
The direction-averaged optimum over all node sets is non-decreasing in $Q$ (restrict to the two faces of a new coordinate). Is $\lim_{Q\to\infty}\beta^*(Q,R)$ of order $1/R$, so that the slab is asymptotically optimal, or exponentially small in $R$?
\end{problem}

\begin{problem}[structure of fully good regions]
Is a connected set of vertices all of whose edges are good the domain of a single global formula, as the whole cube and every entirely good subcube of dimension $<R$ are? If so, each such region is a set on which parity agrees with a polynomial of degree $<R$, of measure at most about $R\sqrt{2/(\pi Q)}$ by the sign-error bound \cite{abfr}, and an isoperimetric inequality for such sets would give the touched-set bound.
\end{problem}

\begin{problem}[linear polynomials]
For $R=2$ the good edges form a $C_4$-free subgraph of the cube, so asymptotically at least $1-0.6068$ of the edges are bad \cite{bhll}, and two vertices with no bad edges are never at distance $2$. Is the optimal bad fraction $\frac12-o(1)$ for generic nodes, and do special nodes (such as $t/5$ at $Q=5$) help asymptotically?
\end{problem}

\begin{problem}[realizing Tur\'an constructions]
The vertex sets of density $2^{-R/2+o(R)}$ meeting every $R$-subcube \cite{alon,eil} are built layer by layer. Is the set of edges at such a vertex set the bad set of some family? A positive answer with $R\le Q^{1/2}$ would refute the touched-set bound in the strongest way, and one for a single moderate $R$ would beat the slab there.
\end{problem}

\section{Formal verification}\label{sec:lean}

The following statements are proved in the Lake project \texttt{EthConsequences} (Lean 4.34.1, Mathlib at a pinned revision). The build passes with \texttt{lake --wfail build}, the only axioms used are \texttt{Classical.choice}, \texttt{Quot.sound} and \texttt{propext}, and the certificate records file hashes and all 47 checked declarations.

\begin{center}\small
\begin{tabular}{>{\raggedright}p{0.34\textwidth}>{\raggedright\arraybackslash}p{0.58\textwidth}}
\toprule
Result & Lean declarations\\
\midrule
Theorem~\ref{thm:barrier}, Corollary~\ref{cor:binomial} & \texttt{RankBarrier.barrier}, \texttt{barrier\_threshold}, \texttt{count\_le\_count\_add\_hammingDist}\\
Proposition~\ref{prop:subcube} (on the full $d$-cube) & \texttt{Subcube.no\_good\_subcube}, \texttt{card\_good\_directions\_lt}, \texttt{eval\_zero\_eq\_sum}\\
Theorem~\ref{thm:local}, necessity & \texttt{Subcube.subcube\_identity}\\
Lemmas~\ref{lem:global}, \ref{lem:symmetric}, Theorem~\ref{thm:slab} & \texttt{SlabDesign.G\_edge}, \texttt{natDegree\_G\_le}, \texttt{G\_eval\_zero}, \texttt{exists\_gamma}, \texttt{slab\_edge\_good}, \texttt{slab\_eval\_zero}, \texttt{slab\_natDegree}, \texttt{slab\_design\_bound}\\
Theorem~\ref{thm:general} & \texttt{PointBarrier.point\_dependent\_barrier}, \texttt{label\_barrier}, \texttt{label\_barrier\_explicit}\\
Theorem~\ref{thm:stress}, soundness direction & \texttt{Stress.stress\_total\_eq\_zero}, \texttt{moment\_collapse}, \texttt{sum\_edgeW}\\
Theorem~\ref{thm:product} & \texttt{Stress.product\_edge\_good}, \texttt{natDegree\_product\_le}, \texttt{product\_eval\_zero}\\
\bottomrule
\end{tabular}
\end{center}

Not formalized: Proposition~\ref{prop:reduction} (checked on actual matrices), the sufficiency half of Theorem~\ref{thm:local} (checked exactly on small cases), the binomial estimates in Corollaries~\ref{cor:degree} and~\ref{cor:helps}, the single-matrix bound in Corollary~\ref{cor:helps}, the completeness direction of Theorem~\ref{thm:stress}, Corollary~\ref{cor:weak}, Propositions~\ref{prop:criterion} and~\ref{prop:forces} and Corollary~\ref{cor:criterion}, and the manuscript's own analysis, on which the last three are conditional.

\section{Status, scope and novelty}\label{sec:status}

\begin{itemize}[leftmargin=2em]
\item \textbf{Scope.} The barriers cover designs built on the label space $([r]\times[r])^Q$ with the group rule $a_w(t)=d_u(t)$: Theorem~\ref{thm:barrier} the template, Theorem~\ref{thm:general} all node-dependent splittings. They say nothing about algorithms of other shapes, nor about the manuscript's Appendix~D, whose square root comes from a partition count of $2$-adic modes.
\item \textbf{Dependence on the main theorem.} None of the theorems depends on it. Propositions~\ref{prop:criterion} and~\ref{prop:forces} and Corollary~\ref{cor:criterion} are conditional on the manuscript's cost analysis.
\item \textbf{Inspection.} We read the manuscript's Sections~1 and~2, the statements of Section~3, Sections~4.1 to~4.3 and Appendices~C and~D in full from the openai/math repository at commit \texttt{7ef5c0cc5}, with the numbering of \texttt{paper.pdf}; this is not a proof audit, and the upstream Lean development was not rebuilt.
\item \textbf{Review.} A separate review pass re-derived the main statements, checked the Lean statements against the claims and the small cases independently, and found the template issue described in the introduction, the dependence of the improvement threshold on the row count, and several smaller errors, all corrected here.
\item \textbf{Novelty.} A web-based literature check on 29 October 2026 found the expositions \cite{odonnell,kintali}, which discuss no barrier and no designs beyond the template, and the closest combinatorics in hypercube Tur\'an problems \cite{aks,eil,alon,bhll} and in the interpolation and probabilistic degree of parity \cite{abfr,alman,stv}. The check was not systematic, so novelty is \textbf{not verified}. A reported independent refutation of ETH by another group, known to us only from a press report, has not been inspected.
\end{itemize}

\begin{thebibliography}{99}\small
\bibitem{oai378} OpenAI, \emph{The Exponential Time Hypothesis is False}, openai/math, family 378, manuscript revised 25 October 2026, \url{https://github.com/openai/math}.
\bibitem{oai378lean} openai/math, Lean scope note for family 378, \texttt{lean/docs/378.md}.
\bibitem{odonnell} R. O'Donnell, \emph{Some exposition of ``The Exponential Time Hypothesis Is False''}, 28 October 2026, \url{https://www.cs.cmu.edu/~odonnell/some_eth_false_exposition.pdf}.
\bibitem{kintali} S. Kintali, \emph{A simplified proof of ETH is false}, 26 October 2026, \url{https://shivakintali.github.io/papers/ETH.pdf}.
\bibitem{abfr} J. Aspnes, R. Beigel, M. Furst, S. Rudich, The expressive power of voting polynomials, \emph{Combinatorica} 14 (1994).
\bibitem{aks} N. Alon, A. Krech, T. Szab\'o, Tur\'an's theorem in the hypercube, \emph{SIAM J. Discrete Math.} 21 (2007).
\bibitem{eil} D. Ellis, M.-R. Ivan, I. Leader, Tur\'an densities for daisies and hypercubes, arXiv:2401.16289.
\bibitem{alon} N. Alon, Erasure list-decodable codes and Tur\'an hypercube problems, \emph{Finite Fields Appl.} 100 (2024), arXiv:2403.16319.
\bibitem{bhll} J. Balogh, P. Hu, B. Lidick\'y, H. Liu, Upper bounds on the size of 4- and 6-cycle-free subgraphs of the hypercube, arXiv:1201.0209.
\bibitem{alman} J. Alman, R. Williams, Probabilistic rank and matrix rigidity, STOC 2017, arXiv:1611.05558.
\bibitem{stv} S. Srinivasan, U. Tripathi, S. Venkitesh, On the probabilistic degrees of symmetric Boolean functions, arXiv:1910.02465.
\end{thebibliography}

\end{document}
